\subsection{What is an algorithm?}

An \textbf{algorithm} is a well-defined computational procedure that takes one or more
values as \textbf{input} and produces one or more values as
\textbf{output}: it is a sequence of steps transforming input into output.

We distinguish an algorithm from a \textbf{data structure}, which is a method for
storing and organizing information to facilitate its access and modification.

% ─────────────────────────────────────────────────────────────────────────────
\subsection{Introductory example: arithmetic sum}

\textbf{Problem:} calculate $\displaystyle S = \sum_{i=1}^{n} i$.

\begin{center}
\begin{tikzpicture}[>=Stealth]
  % Naive box
  \node[draw=keyred, fill=red!8, rounded corners, text width=5.5cm, align=center,
        inner sep=8pt] (naive) {
    \textbf{\textcolor{keyred}{Naive approach}}\\[4pt]
    Loop: add $1 + 2 + \cdots + n$\\[4pt]
    $n$ additions \;\;$\Rightarrow$\;\; $\Theta(n)$
  };
  % Smart box
  \node[draw=keygreen, fill=green!8, rounded corners, text width=5.5cm, align=center,
        inner sep=8pt, right=1.5cm of naive] (smart) {
    \textbf{\textcolor{keygreen}{Gauss's approach}}\\[4pt]
    Closed-form formula: $\dfrac{n(n+1)}{2}$\\[4pt]
    3 operations \;\;$\Rightarrow$\;\; $\Theta(1)$
  };
  \draw[->, thick, gray] (naive) -- node[above, font=\small\itshape]{vs.} (smart);
\end{tikzpicture}
\end{center}

% ─────────────────────────────────────────────────────────────────────────────
\subsection{Why does efficiency matter?}

A \emph{slow} computer running a good algorithm will always eventually beat a
\emph{fast} computer running a bad algorithm, as soon as the input size $n$
becomes large enough.

\begin{center}
\begin{tikzpicture}
  \begin{scope}
    \draw[->] (0,0) -- (5.5,0) node[right] {$n$};
    \draw[->] (0,0) -- (0,4.2) node[above] {time};
    % O(n^2)
    \draw[keyred, thick, domain=0:2.0, samples=60]
          plot (\x, {0.4*\x*\x}) node[right] {$O(n^2)$};
    % O(n log n)
    \draw[keyblue, thick, domain=0:5, samples=100]
          plot (\x, {0.35*\x*(ln(\x+1)/ln(2)+0.1)}) node[right] {$O(n\log n)$};
    % O(n)
    \draw[keygreen, thick, domain=0:5, samples=60]
          plot (\x, {0.4*\x}) node[right] {$O(n)$};
    % O(1)
    \draw[orange, thick] (0,0.3) -- (5,0.3) node[right] {$O(1)$};
  \end{scope}
\end{tikzpicture}
\end{center}

\begin{tcolorbox}[colback=lightyellow, colframe=orange!70!black, fonttitle=\bfseries,
                  title={Conclusion}]
The \textbf{asymptotic performance} (behavior as $n \to \infty$) is more
decisive than the multiplicative constants related to the hardware.
\end{tcolorbox}